% Solutions of Book 4, chapter 2 (cpp-copy-move).
\section*{Chapter \ref{ch:b2:cpp-copy-move} --- Copying, Moving and Smart Pointers}

\begin{solution}{exo:b2:cpp-copy-move:1}
A: \code{a} is an \omterm{def:b2:cpp-copy-move:rvalue}{lvalue}, copy. B: \code{std::move(a)} is an xvalue, move. C: assignment from the \omterm{def:b2:cpp-copy-move:rvalue}{lvalue} \code{c},
copy assignment. D: the \omterm{def:b2:cpp-copy-move:rvalue}{lvalue} \code{b} initialises the parameter \code{t} by copy; \code{return t;} cannot be elided
(a parameter is not built in the return slot) and moves. E: the \omterm{def:b2:cpp-copy-move:rvalue}{lvalue} \code{d} is copied into the vector. F: the
prvalue \code{Tracked(5)} binds to \code{push_back(T &&)} and is moved in; thanks to \code{reserve(4)} no regrowth
carries anything. Total: copies A, C, D, E and moves B, D, F.
\omoutput{code/bachelor-2/02-cpp-copy-move/outputs/exo_count.txt}{Four copies, three moves.}
\end{solution}

\begin{solution}{exo:b2:cpp-copy-move:2}
\code{o} is \code{*this}, so \code{delete[] data_} frees the very text that \code{strlen(o.data_)} then reads: the
report says so, a read of freed memory inside \code{Text::operator=} on line 15, called from line 29. First fix:
begin with \code{if (this == &o) return *this;}. It cures self-assignment, but if \code{new} then throws, \code{data_}
still holds the freed pointer and the \omterm{def:b2:cpp-classes-raii:constructor}{destructor} will free it again: the operation does not even give the \omterm{def:b2:cpp-classes-raii:guarantee}{basic
guarantee}. Second fix, by \cref{met:b2:cpp-copy-move:copy-and-swap}: \code{Text tmp(o); std::swap(data_, tmp.data_);
return *this;}. The copy is made before anything is freed, so self-assignment works without a test, and by
\cref{prop:b2:cpp-copy-move:copy-and-swap} a throwing \code{new} leaves the target unchanged: only this fix is
strong.
\end{solution}

\begin{solution}{exo:b2:cpp-copy-move:3}
\omcode{code/bachelor-2/02-cpp-copy-move/exercises/buffer5.hpp}{7}{32}{\code{Buffer} by the \omterm{def:b2:cpp-copy-move:rules}{rule of five}}
\code{std::exchange(x, v)} sets \code{x} to \code{v} and returns its old value, which writes ``take and empty'' in one
expression. The parameter \code{o} of \code{operator=} is a new object built by the caller: by the \omterm{def:b2:cpp-copy-move:copy}{copy constructor} if
the argument is an \omterm{def:b2:cpp-copy-move:rvalue}{lvalue}, by the \omterm{def:b2:cpp-copy-move:move}{move constructor} (or directly, by elision) if it is an \omterm{def:b2:cpp-copy-move:rvalue}{rvalue}. Either way the body
only swaps, and \code{o} takes the old block away with it. One function therefore provides both assignments; the
price is a \omterm{def:b2:cpp-copy-move:move}{move constructor} call more than a dedicated move assignment would make. Tested in
\code{tests/buffer_test.cpp}.
\end{solution}

\begin{solution}{exo:b2:cpp-copy-move:4}
With $n = 2^m$ elements, the regrowths happen when the array is full at capacities $1, 2, 4, \dots, 2^{m-1}$, which
carry $2^m - 1 = n - 1$ elements. With $n = 2^m + 1$ there is one more regrowth, at $2^m$, and the total is
$2^{m+1} - 1 = 2n - 3$. Both are less than $2(n-1)$, the bound of \cref{prop:b2:cpp-copy-move:growth-moves} for
$r = 2$. The run: $n = 1000$ lies between $2^9 + 1$ and $2^{10}$, so the regrowths are those of $2^{10}$, 1023
transfers; $n = 1025 = 2^{10} + 1$ gives $2 \cdot 1025 - 3 = 2047$. Both match \cref{out:b2:cpp-copy-move:growth}.
\end{solution}

\begin{solution}{exo:b2:cpp-copy-move:5}
\omoutput{code/bachelor-2/02-cpp-copy-move/outputs/exo_which.txt}{The ten lines.}
1: default \omterm{def:b2:cpp-classes-raii:constructor}{constructor}. 2: \code{a} is an \omterm{def:b2:cpp-copy-move:rvalue}{lvalue}, \omterm{def:b2:cpp-copy-move:copy}{copy constructor}. 3: an xvalue, \omterm{def:b2:cpp-copy-move:move}{move constructor}. 4: assignment
from an \omterm{def:b2:cpp-copy-move:rvalue}{lvalue}, copy assignment. 5: \code{make()} builds its \code{W()} (\code{default}) in a temporary, by guaranteed
elision, and the temporary, an \omterm{def:b2:cpp-copy-move:rvalue}{rvalue}, is move-assigned to \code{a}. 6: guaranteed elision builds the result of
\code{make} directly in \code{d}: \code{default} only. 7: \code{named} returns a named local; this compiler applied
the named return value optimisation, so only \code{default} is printed. That is the statement that depends on the
compiler: without the optimisation it prints \code{default} then \code{move constructor}. 8: \code{f} is built
(\code{default}); \code{std::move(f)} has type \code{const W &&}, which cannot bind to \code{W &&} but binds to
\code{const W &}: the \omterm{def:b2:cpp-copy-move:copy}{copy constructor} runs.
\end{solution}

\begin{solution}{exo:b2:cpp-copy-move:6}
The first \code{push_back(std::move(jobs))} move-constructs the day's list in \code{history}, and a vector's \omterm{def:b2:cpp-copy-move:move}{move
constructor} leaves its source empty; from day 2 on, the empty \code{jobs} is moved, which gives empty lists. This is
not undefined behaviour: a moved-from vector is valid, and calling \code{size} on it is allowed; the program does
exactly what it says. The fix keeps \code{jobs} and copies it: \code{history.push_back(jobs);}, which is the line that
copies (the three strings are copied, three times in all). If \code{jobs} is not needed after the loop, the last day
may still move it.
\end{solution}

\begin{solution}{exo:b2:cpp-copy-move:7}
\omcode{code/bachelor-2/02-cpp-copy-move/exercises/unique_ptr.hpp}{5}{32}{\code{UniquePtr}}
The move assignment tests \code{this != &o}: without the test, a self-move would \code{reset} with the pointer it
just took from itself and leave a null pointer behind, after deleting nothing; with \code{reset} written as above, the
order ``store the new pointer, then delete the old'' also keeps the class correct when the deleted object's
\omterm{def:b2:cpp-classes-raii:constructor}{destructor} reaches back to its owner. \code{tests/exercises_test.cpp} counts live objects: a hundred
\code{UniquePtr} pushed into a vector, which moves them at each regrowth, are freed exactly once each when the vector
goes, and a moved-from pointer is null.
\end{solution}

\begin{solution}{exo:b2:cpp-copy-move:8}
With \code{Buffer &operator=(Buffer o) noexcept}, the copy that may throw is the initialisation of the parameter
\code{o}, which happens at the call, before the body of \code{operator=} starts. If it throws, \code{operator=} is
never entered and no statement has touched the target: it keeps its value, which is the strong guarantee. If it
succeeds, the body only swaps pointers and numbers, which cannot throw, and the \omterm{def:b2:cpp-classes-raii:constructor}{destructor} of \code{o} frees the old
block, which does not throw either. For an \omterm{def:b2:cpp-copy-move:rvalue}{rvalue} argument, \code{o} is built by the \code{noexcept} \omterm{def:b2:cpp-copy-move:move}{move
constructor}, so the whole assignment cannot throw: it gives the no-throw guarantee, which implies the strong one.
Self-assignment works as before, since \code{o} is a separate copy.
\end{solution}

\begin{solution}{exo:b2:cpp-copy-move:9}
\omcode{code/bachelor-2/02-cpp-copy-move/exercises/tree.hpp}{7}{31}{A tree that owns downward and observes upward}
A node is owned by the vector of its parent, so it is destroyed no later than its parent: whenever a node exists,
its parent exists, and the raw pointer never dangles; \code{move_child} keeps this true by updating it. Had each child
held a \code{std::shared_ptr} to its parent, with the parent holding its children by \code{shared_ptr}, each
parent--child pair would be a cycle of owners and the whole tree would leak, as in
\cref{out:b2:cpp-copy-move:shared-cycle}. The moved subtree keeps its address: only the \code{unique_ptr} that owns
it moves from one vector to the other, and the test checks that \code{&mill} is unchanged.
\end{solution}

\begin{solution}{exo:b2:cpp-copy-move:10}
\omcodes{code/bachelor-2/02-cpp-copy-move/exercises/small_string.hpp}{11-40,46-55}{\code{SmallString}}
A small string's pointer points into the object's own array \code{buf_}. Stealing it would leave the target pointing
into the source, which is about to die; the \omterm{def:b2:cpp-copy-move:move}{move constructor} must copy the characters into its own array and point
to that. So moving a small string costs what copying it costs, at most 16 bytes, while moving a long one takes its
block. The move assignment destroys the target and rebuilds it in place with the \omterm{def:b2:cpp-copy-move:move}{move constructor}, which is correct
because that \omterm{def:b2:cpp-classes-raii:constructor}{constructor} cannot throw (otherwise a throw would leave a destroyed object behind); a swap-based
version would be as good. \code{tests/exercises_test.cpp} covers copies and moves between small and long strings,
self-assignment and self-move, and a vector of fifty strings that grows by moving them.
\end{solution}

\begin{solution}{exo:b2:cpp-copy-move:11}
(a) For $r = 2$ and $n = 2^m + 1$, \cref{exo:b2:cpp-copy-move:4} found $2n - 3$ transfers, and
$(2n-3)/(2n-2) \to 1$ as $m \to \infty$: no constant $\alpha < 1$ can multiply the bound. (b) The proof gives
$3(n-1)$ transfers for any growth by a factor at least $1.5$. The rounded rule does not satisfy that hypothesis
($3 \to 4$ is a factor $1.33$, $9 \to 13$ is $1.44$), so the proposition does not apply as stated; counting directly,
the array is full and regrows at $1, 2, 3, 4, 6, 9, 13, 19, 28, 42, 63, 94$, which carry 284 elements for $n = 100$,
still under $3 \cdot 99 = 297$. (c) With capacities $d, 2d, 3d, \dots$, the regrowths up to $n$ carry $d + 2d + \dots
+ kd = d\,k(k+1)/2$ elements with $k = \lfloor (n-1)/d \rfloor$, about $n^2/(2d)$: quadratic in $n$, so the amortised
cost of an append grows like $n/(2d)$ instead of staying constant. Only a multiplicative growth gives constant
amortised cost.
\end{solution}

\begin{solution}{exo:b2:cpp-copy-move:12}
\code{station count 3}: \code{s} and the \code{source} of each display. When \code{main} returns, \code{s} goes and
the station's count falls to 2; each display is owned once, by the station's vector. The four objects own each
other in two cycles (station and hall, station and office), nothing outside reaches them, and LeakSanitizer reports
four blocks. The fix makes \code{Display::source} a \code{std::weak_ptr<Station>}. The other edge cannot be the weak
one: the station's vector is the only owner of each display, so with weak edges there the displays would die as soon
as \code{watch} returned.
\omcode{code/bachelor-2/02-cpp-copy-move/exercises/observers_fixed.hpp}{10}{29}{The fixed classes (they also count
the live objects)}
Proof that everything is freed. In the fixed version the station's strong count counts only \omterm{def:b2:cpp-copy-move:shared-ptr}{shared pointers} held
outside the station and its displays, since the displays hold weak ones. When the last of them goes the count reaches
0, so the station is destroyed; its members are destroyed, among them the vector, which destroys its \omterm{def:b2:cpp-copy-move:shared-ptr}{shared pointers}
to the displays. Each display had exactly one strong owner, that vector, so each count reaches 0 and each display is
destroyed; its \code{weak_ptr} member does not keep anything alive. This holds for any number of displays, and the
test checks that the count of live objects returns to 0.
\end{solution}

\begin{solution}{pb:b2:cpp-copy-move:1}
\omoutput{code/bachelor-2/02-cpp-copy-move/outputs/pipeline.txt}{The four reports.}

\textbf{1.} $4000 \times 3000 = 12\,000\,000$ bytes, 12\,MB, allocated and copied.

\textbf{2.} \code{Image out = im;} in \code{brighten1} and the same line in \code{darken1}: 2 copies, 0 moves.

\textbf{3.} The compiler built each \code{out} directly in the caller's object (named return value optimisation).
Without it each return would move: 2 copies and 2 moves.

\textbf{4.} When \code{photo}, an \omterm{def:b2:cpp-copy-move:rvalue}{lvalue}, initialises the parameter \code{im} of \code{brighten2}. The edit changes
\code{im}; if it shared \code{photo}'s pixels, \code{photo} would change too, and \omterm{def:b2:cpp-copy-move:value-semantics}{value semantics} forbid that.

\textbf{5.} \code{return im;} in \code{brighten2} and in \code{darken2}. A parameter is built by the caller before the
call, in a place chosen before the return slot is known, so it cannot be built in the return slot; the language then
treats the returned name as an \omterm{def:b2:cpp-copy-move:rvalue}{rvalue} and moves.

\textbf{6.} \code{brighten2(photo)} is a prvalue of type \code{Image}; it initialises the parameter of \code{darken2}
directly (guaranteed elision).

\textbf{7.} 0 copies, 3 moves: the parameter of \code{brighten2} is now move-constructed. \code{photo} is
moved-from: the program may destroy it or assign it a new value, and should not read its pixels.

\textbf{8.} First push: capacity 0 to 1, \code{a} copied (1 copy). Second: capacity 1 to 2, the one element moved
over and \code{b} moved in (2 moves). Third: 2 to 4, two elements moved over and \code{c} moved in (3 moves). Total
1 copy and 5 moves, as printed.

\textbf{9.} Four copies in all, 48\,MB. Without \code{noexcept}, the three transfers of regrowth become copies (the
other moves are unchanged, since only the vector checks \code{noexcept}): seven copies, 84\,MB.

\textbf{10.} No regrowth: 1 copy and 2 moves.

\textbf{11.} All five: the member-wise copy copies the pixel vector (a \omterm{def:b2:cpp-copy-move:deep}{deep copy}, the same work) and the member-wise
move moves it, and the compiler's moves are \code{noexcept} because the vector's are. Only the counting is lost.

\textbf{12.} If the function needs its own copy of the argument to change it, take it by value: a caller who still
needs its object pays the one copy it needs, and a caller with a temporary or a \code{std::move} pays nothing. If the
function only reads, take \code{const Image &}. Version 1 copies even when the caller had a temporary.

\textbf{13.} A \code{unique_ptr} is the sole owner, and here a photo has one owner per album, none of which knows
whether it will be the last.

\textbf{14.} Album 1 $\to$ photo and album 2 $\to$ photo (strong), photo $\to$ album 1 and photo $\to$ album 2
(strong). Once the library's pointers are gone: photo strong 2, each album strong 1.

\textbf{15.} All three objects: each has a positive count and none is reachable. At least the photo's 12\,000\,000
pixel bytes, plus the three objects and their control blocks.

\textbf{16.} The edges from the photo to its albums become weak: a photo observes the albums that contain it. While
the library holds the albums: each album strong 1 and weak 1, the photo strong 2. When the library drops an album,
its count reaches 0 and it is freed, which drops the photo to strong 1; when the second goes, the photo is freed too.

\textbf{17.} With C++20, \code{std::erase_if(albums, gone);} where \code{gone} is the lambda
\code{[](const std::weak_ptr<Album> &w) {return w.expired();}}; or a loop that keeps the entries for which
\code{w.lock()} is not null.

\textbf{18.} The drawing code holds a \code{std::shared_ptr<Photo>} for the whole drawing, obtained by copying
another or by \code{lock()} on a \code{weak_ptr}: it is then one of the owners, and the count cannot reach 0 before it
lets go. (Updating the count from several threads at once is the subject of \cref{ch:b2:threads-and-synchronisation}.)

\textbf{19.} About 39\,ns for $10^3$ elements, a few hundredths of a nanosecond per element; about 5\,ms for $10^7$,
half a nanosecond per element. The small vectors stay in the caches, from which the processor copies many bytes per
cycle; the large ones come from main memory, whose rate is about ten times lower.

\textbf{20.} Assuming a byte costs what it costs inside a large vector of \code{double}, about $0.5/8$\,ns, a copy of
12\,MB takes of the order of $0.75$\,ms.

\textbf{21.} Version 1 makes 2 copies per photo: $50\,000 \times 2 \times 0.75$\,ms, of the order of a minute and a
quarter per night. Version 2 with \code{std::move} makes none and saves all of it.

\textbf{22.} A move of a vector reads the three pointers of the source (start, end of the elements, end of the
capacity), writes them into the target, and sets the source's to null; it never touches the elements, so its time
does not depend on $n$.

\textbf{23.} With $n = 1000$ and doubling, fewer than $2 \cdot 999 = 1998$ transfers (in fact 1023); as copies of
12\,MB each, at most about 24\,GB copied (in fact about 12\,GB). With \code{noexcept} moves, no pixel is copied.

\textbf{24.} \emph{Transfers of a growing array: fewer than $\frac{r}{r-1}(n-1)$ for a growth factor $r$, that is
$2(n-1)$ for doubling, and zero copies with \code{noexcept} moves}: each transfer then costs a few pointer copies
instead of a copy of the element's data, 12\,MB per image here.
\end{solution}
